\documentclass{beamer}

\usetheme{Madrid}
\usecolortheme{default}

\usepackage{graphicx}
\usepackage{amsmath}
\usepackage{amsfonts}
\usepackage{booktabs}
\usepackage{array}
\usepackage{caption}

\title[Vision-Language-Action + Diffusion Switching]{Vision-Language-Action Model and Diffusion Policy Switching Enables Dexterous Control of an Anthropomorphic Hand}
\author[Pan, Junge, Hughes]{Cheng Pan \and Kai Junge \and Josie Hughes}
\institute[EPFL]{Faculty of Mechanical Engineering\\Swiss Federal Institute of Technology Lausanne (EPFL)\\Lausanne, Switzerland}
\date{2024}

\begin{document}

%------------------------------------------------
\begin{frame}
  \titlepage
\end{frame}

%------------------------------------------------
\begin{frame}
  \frametitle{Citation}
  \small
  Pan, C., Junge, K., \& Hughes, J. (2024). \textit{Vision-Language-Action Model and Diffusion Policy Switching Enables Dexterous Control of an Anthropomorphic Hand}.\\[0.4cm]

  \textbf{Type:} Robotics / Manipulation, Vision-Language-Action and Diffusion Policies\\[0.2cm]
  \textbf{Affiliation:} Faculty of Mechanical Engineering, EPFL, Switzerland\\[0.2cm]
  \textbf{Preprint:} Project page: \texttt{https://vla-diffu-switch.github.io/}\\[0.2cm]

  \textbf{Keywords:}
  \begin{itemize}
    \item Vision-Language-Action (VLA) models
    \item Diffusion policy for robot control
    \item Dexterous manipulation, anthropomorphic hand
    \item Hybrid high-level / low-level control
  \end{itemize}
\end{frame}

%------------------------------------------------
\begin{frame}
  \frametitle{Research Area and Background (1/2)}
  \small
  \begin{itemize}
    \item Dexterous manipulation requires:
      \begin{itemize}
        \item Advanced task planning
        \item Wrist trajectory generation
        \item Contextual grasp selection
        \item Precise multi-finger control
      \end{itemize}
    \item Large Vision-Language-Action (VLA) models:
      \begin{itemize}
        \item Translate language commands to robot actions
        \item Trained on large-scale manipulation datasets (e.g., Open-X-Embodiment, DROID, BridgeData v2)
      \end{itemize}
    \item Examples: RT-1, RT-2, Octo, openVLA
  \end{itemize}
\end{frame}

%------------------------------------------------
\begin{frame}
  \frametitle{Research Area and Background (2/2)}
  \small
  \begin{itemize}
    \item Existing VLA deployments:
      \begin{itemize}
        \item Primarily on 1 DoF pinch grippers
        \item Limited evidence on multi-fingered, anthropomorphic hands
      \end{itemize}
    \item Diffusion policy models:
      \begin{itemize}
        \item Well-suited for robot motion learning and planning
        \item Can generate smooth, precise motor commands
        \item Capture contextual information from sensor inputs
      \end{itemize}
    \item Prior diffusion work includes multi-fingered hands, but typically task-specific
  \end{itemize}
\end{frame}

%------------------------------------------------
\begin{frame}
  \frametitle{Goal of the Research}
  \small
  \begin{itemize}
    \item Develop a \textbf{hybrid control framework} combining:
      \begin{itemize}
        \item Fine-tuned openVLA model for \textbf{language-conditioned high-level planning}
        \item Diffusion policy for \textbf{low-level dexterous grasping control}
      \end{itemize}
    \item Enable \textbf{autonomous pick-and-place} with:
      \begin{itemize}
        \item Language-specified target object and placement location
        \item Multi-fingered anthropomorphic hand (ADAPT Hand 2)
      \end{itemize}
    \item Demonstrate \textbf{event-based switching} between VLA and diffusion policies
    \item Achieve higher success rates than VLA-only control
  \end{itemize}
\end{frame}

%------------------------------------------------
\begin{frame}
  \frametitle{State-of-the-Art and Limitations}
  \small
  \begin{itemize}
    \item \textbf{VLAs (e.g., RT-1, RT-2, Octo, openVLA):}
      \begin{itemize}
        \item Strength: generalizable language-to-action mapping
        \item Limitation: demonstrated mainly on simple grippers
        \item Unclear precision for complex anthropomorphic hands
      \end{itemize}
    \item \textbf{Visual servoing:}
      \begin{itemize}
        \item Effective for approaching and grasping
        \item Limited use of contextual/environmental information for grasp strategy
      \end{itemize}
    \item \textbf{Diffusion policies:}
      \begin{itemize}
        \item Strength: precise, smooth control; multi-modal behavior
        \item Limitation: task- and object-specific training; limited global planning
      \end{itemize}
  \end{itemize}
\end{frame}

%------------------------------------------------
\begin{frame}
  \frametitle{Research Motivation}
  \small
  \begin{itemize}
    \item VLAs:
      \begin{itemize}
        \item Good at high-level planning from language
        \item Struggle with precise, contact-rich dexterous manipulation
      \end{itemize}
    \item Diffusion policies:
      \begin{itemize}
        \item Good at local, precise control and exploiting context
        \item Lack global task planning and language grounding
      \end{itemize}
    \item \textbf{Key idea:} Combine VLA and diffusion to exploit complementary strengths
      \begin{itemize}
        \item VLA: move hand near target from language
        \item Diffusion: execute robust grasp and release
      \end{itemize}
    \item Target platform: compliant anthropomorphic ADAPT Hand 2 (13 DoF)
  \end{itemize}
\end{frame}

%------------------------------------------------
\begin{frame}
  \frametitle{Problem Definition}
  \small
  \begin{itemize}
    \item Task: \textbf{Autonomous pick-and-place} with:
      \begin{itemize}
        \item Language command: ``pick [object] and place on [plate color]''
        \item Multiple objects and placement locations
      \end{itemize}
    \item Robot:
      \begin{itemize}
        \item ADAPT Hand 2 (13 DoF) mounted on UR5 arm
        \item Two RGB cameras: static (cam1) and wrist-mounted (cam2)
      \end{itemize}
    \item Requirements:
      \begin{itemize}
        \item Language-conditioned high-level motion
        \item Precise dexterous grasping and release
        \item Automatic switching between controllers
      \end{itemize}
    \item Baseline: VLA-only control for full pick-and-place
  \end{itemize}
\end{frame}

%------------------------------------------------
\begin{frame}
  \frametitle{Proposed Hybrid Method: Overview}
  \small
  \begin{columns}
    \begin{column}{0.55\textwidth}
      \begin{itemize}
        \item \textbf{VLA (openVLA, fine-tuned):}
          \begin{itemize}
            \item Input: language command + images (cam1+cam2)
            \item Output: 6D end-effector pose + scalar grasp/event signal
          \end{itemize}
        \item \textbf{Diffusion policy:}
          \begin{itemize}
            \item Input: images (cam1, cam2) + state
            \item Output: multi-dimensional hand joint commands
          \end{itemize}
        \item \textbf{Event signal $\sigma$:}
          \begin{itemize}
            \item Shared scalar signal to mark key task events
            \item Drives automatic switching between policies
          \end{itemize}
      \end{itemize}
    \end{column}
    \begin{column}{0.45\textwidth}
      \centering
      \includegraphics[width=0.7\textwidth,keepaspectratio]{extracted/figure_1.png}
      \captionof{figure}{Combined VLA and diffusion policy approach for dexterous manipulation (Fig.~1).}
    \end{column}
  \end{columns}
\end{frame}

%------------------------------------------------
\begin{frame}
  \frametitle{VLA \& Diffusion Switching Framework}
  \small
  \begin{columns}
    \begin{column}{0.55\textwidth}
      \begin{itemize}
        \item Use a \textbf{common event signal} $\sigma$ to:
          \begin{itemize}
            \item Detect proximity to object / plate (VLA)
            \item Detect successful grasp (diffusion)
          \end{itemize}
        \item Execution sequence:
          \begin{itemize}
            \item VLA moves hand near object; $\sigma$ rises
            \item Switch to diffusion for grasping; $\sigma$ rises at successful lift
            \item Switch back to VLA for transport and placement
          \end{itemize}
        \item Object-specific diffusion policies:
          \begin{itemize}
            \item Selected via lookup table from language input
          \end{itemize}
      \end{itemize}
    \end{column}
    \begin{column}{0.45\textwidth}
      \centering
      \includegraphics[width=0.7\textwidth,keepaspectratio]{extracted/figure_2.png}
      \captionof{figure}{Switching between VLA and diffusion using event signal $\sigma$ (Fig.~2).}
    \end{column}
  \end{columns}
\end{frame}

%------------------------------------------------
\begin{frame}
  \frametitle{Robotic Platform: ADAPT Hand 2}
  \small
  \begin{columns}
    \begin{column}{0.55\textwidth}
      \begin{itemize}
        \item Anthropomorphic hand with 13 motors
        \item Series elastic actuation at base joints:
          \begin{itemize}
            \item Compliance for safe interaction
            \item Physical filtering of noisy control signals
          \end{itemize}
        \item Anatomical kinematic design and continuous soft skin
        \item Mounted on UR5 arm for 6-DoF positioning
      \end{itemize}
    \end{column}
    \begin{column}{0.45\textwidth}
      \centering
      \includegraphics[width=0.7\textwidth,keepaspectratio]{extracted/figure_3.png}
      \captionof{figure}{ADAPT Hand 2 with compliant joints and continuous skin (Fig.~3).}
    \end{column}
  \end{columns}
\end{frame}

%------------------------------------------------
\begin{frame}
  \frametitle{Experimental Setup and Teleoperation}
  \small
  \begin{columns}
    \begin{column}{0.55\textwidth}
      \begin{itemize}
        \item Teleoperation for data collection:
          \begin{itemize}
            \item Apple Vision Pro hand tracking
            \item Joint-to-joint mapping to UR5 + ADAPT Hand 2
          \end{itemize}
        \item Sensors:
          \begin{itemize}
            \item cam1: static world camera
            \item cam2: wrist-mounted camera
          \end{itemize}
        \item Objects:
          \begin{itemize}
            \item Red pepper, tape roll, piece of paper
            \item Two plates (yellow, purple) as placement targets
          \end{itemize}
        \item Inference hardware: NVIDIA RTX 4090, $\approx$5 Hz control
      \end{itemize}
    \end{column}
    \begin{column}{0.45\textwidth}
      \centering
      \includegraphics[width=0.7\textwidth,keepaspectratio]{extracted/figure_4.png}
      \captionof{figure}{Teleoperation setup and test objects/plates (Fig.~4).}
    \end{column}
  \end{columns}
\end{frame}

%------------------------------------------------
\begin{frame}
  \frametitle{Data Collection Strategy}
  \small
  \begin{itemize}
    \item \textbf{VLA dataset (Fig.~5A):}
      \begin{itemize}
        \item Teleoperated full pick-and-place trajectories
        \item Operator records event signal via controller
        \item Only pre- and post-grasp segments used for training
        \item Hand deliberately closed in mid-air when event signal first applied
        \item 20 trials per object-placement combination
      \end{itemize}
    \item \textbf{Diffusion dataset (Fig.~5B):}
      \begin{itemize}
        \item Only grasping segments
        \item Hand initialized $\approx$5 cm above object
        \item 40 trials (tape, paper), 30 trials (pepper)
        \item Includes multiple grasp strategies and failure-recovery attempts
        \item Event signal non-zero only in final 10 steps of successful trials
      \end{itemize}
  \end{itemize}
\end{frame}

%------------------------------------------------
\begin{frame}
  \frametitle{Data Collection Illustration}
  \small
  \centering
  \includegraphics[width=0.7\textwidth,keepaspectratio]{extracted/figure_5.png}
  \captionof{figure}{Data collection for VLA (A) and diffusion policy (B), including event signal annotation (Fig.~5).}
\end{frame}

%------------------------------------------------
\begin{frame}
  \frametitle{Model Training Details}
  \small
  \begin{itemize}
    \item \textbf{VLA (openVLA fine-tuning):}
      \begin{itemize}
        \item Inputs: concatenated images from cam1 (224$\times$144) and cam2 (224$\times$80)
        \item Images vertically concatenated into single frame
        \item Default openVLA fine-tuning settings, batch size 22
        \item No image augmentation (fixed camera pose and lighting)
        \item Trained until action accuracy $>$ 95\% ($\approx$50k steps)
        \item Trained on A100-80GB GPU
      \end{itemize}
    \item \textbf{Diffusion policy:}
      \begin{itemize}
        \item Inputs: images from cam1, cam2 (320$\times$240), random crop to 288$\times$216
        \item CNN-based diffusion policy as in prior work
        \item Trained for 1500 epochs on custom GPU
      \end{itemize}
  \end{itemize}
\end{frame}

%------------------------------------------------
\begin{frame}
  \frametitle{Experimental Design}
  \small
  \begin{itemize}
    \item \textbf{VLA evaluation:}
      \begin{itemize}
        \item Full grasping and placing with VLA-only
        \item Objects: pepper, tape (blue paper excluded due to convergence issues)
        \item 5 repeats per object-plate combination
      \end{itemize}
    \item \textbf{Diffusion evaluation:}
      \begin{itemize}
        \item Success vs. initial hand offset (5, 10, 15 cm)
        \item Multi-modal grasping (slide \& pick vs. direct pick)
        \item Failure recovery behavior
      \end{itemize}
    \item \textbf{Combined VLA + diffusion:}
      \begin{itemize}
        \item Full pick-and-place for all object-plate combinations
        \item 5 trials per combination
        \item Compare success scores to VLA-only baseline
      \end{itemize}
  \end{itemize}
\end{frame}

%------------------------------------------------
\begin{frame}
  \frametitle{Evaluation Metric for Pick-and-Place}
  \small
  \begin{itemize}
    \item Success scored on a 0--1 scale:
      \begin{itemize}
        \item 1.00: Full task success
        \item 0.75: Failed to place on correct plate
        \item 0.50: Placed on incorrect plate
        \item 0.25: Failed to grasp correct object
        \item 0.00: Approached wrong object
      \end{itemize}
    \item Allows more nuanced comparison than binary success
    \item Applied to both VLA-only and VLA+diffusion conditions
  \end{itemize}
\end{frame}

%------------------------------------------------
\begin{frame}
  \frametitle{VLA-Only Performance: Success Rates}
  \small
  \begin{itemize}
    \item VLA fine-tuned on full grasping and placing for pepper and tape
    \item Achieved $>$95\% action accuracy in training
    \item Test success scores (Table I):
  \end{itemize}
  \vspace{0.2cm}
  \centering
  \small
  \begin{tabular}{lcc}
    \toprule
    End Target & Pepper & Tape \\
    \midrule
    Purple plate & 0.45 & 0.20 \\
    Yellow plate & 0.40 & 0.25 \\
    \bottomrule
  \end{tabular}
  \vspace{0.2cm}

  \small
  \begin{itemize}
    \item VLA-only can occasionally grasp pepper
    \item Fails consistently on tape, which requires more precise dexterity
    \item Scores mainly reflect ability to approach correct object
  \end{itemize}
\end{frame}

%------------------------------------------------
\begin{frame}
  \frametitle{VLA Reaching Accuracy}
  \small
  \begin{columns}
    \begin{column}{0.55\textwidth}
      \begin{itemize}
        \item Evaluate offset between hand center and object in x-y plane
        \item Position recorded when VLA event signal triggers switch
        \item Compare:
          \begin{itemize}
            \item VLA with cam1 only
            \item VLA with concatenated cam1+cam2
          \end{itemize}
        \item Results (Fig.~6):
          \begin{itemize}
            \item Combined cameras: typical offset $\leq$ 6 cm
            \item For tape: mean offset reduced from 12 cm to 3 cm
            \item Two cameras help compensate for depth ambiguity
          \end{itemize}
      \end{itemize}
    \end{column}
    \begin{column}{0.45\textwidth}
      \centering
      \includegraphics[width=0.7\textwidth,keepaspectratio]{extracted/figure_6.png}
      \captionof{figure}{x-y offset from target object for different camera setups (Fig.~6).}
    \end{column}
  \end{columns}
\end{frame}

%------------------------------------------------
\begin{frame}
  \frametitle{Diffusion Policy: Multi-Modal Grasping}
  \small
  \begin{columns}
    \begin{column}{0.55\textwidth}
      \begin{itemize}
        \item Objects: tape and blue block
        \item Training includes:
          \begin{itemize}
            \item Slide \& pick trajectories
            \item Direct pick trajectories
            \item Equal number of each type
          \end{itemize}
        \item Test: vary distance from table edge
        \item Results (Fig.~7):
          \begin{itemize}
            \item Closer to edge: higher probability of slide \& pick
            \item Further from edge: higher probability of direct pick
            \item All trials (5 per distance) succeed for both objects
          \end{itemize}
      \end{itemize}
    \end{column}
    \begin{column}{0.45\textwidth}
      \centering
      \includegraphics[width=0.7\textwidth,keepaspectratio]{extracted/figure_7.png}
      \captionof{figure}{Multi-modal grasp selection vs. distance to table edge (Fig.~7).}
    \end{column}
  \end{columns}
\end{frame}

%------------------------------------------------
\begin{frame}
  \frametitle{Diffusion Policy: Sensitivity to Initial Offset}
  \small
  \begin{columns}
    \begin{column}{0.55\textwidth}
      \begin{itemize}
        \item Test grasp success vs. initial hand-object offset
        \item Offsets in x-y plane: 5, 10, 15 cm
        \item Objects: pepper and tape
        \item Results (Fig.~8):
          \begin{itemize}
            \item Success rate drops below 50\% when offset $>$ 15 cm
            \item Indicates diffusion requires hand to start near object
            \item VLA's typical 6 cm error lies within diffusion's tolerance
          \end{itemize}
      \end{itemize}
    \end{column}
    \begin{column}{0.45\textwidth}
      \centering
      \includegraphics[width=0.7\textwidth,keepaspectratio]{extracted/figure_8.png}
      \captionof{figure}{Grasp success vs. initial offset for diffusion policy (Fig.~8).}
    \end{column}
  \end{columns}
\end{frame}

%------------------------------------------------
\begin{frame}
  \frametitle{Diffusion Policy: Failure Recovery}
  \small
  \begin{columns}
    \begin{column}{0.55\textwidth}
      \begin{itemize}
        \item Demonstration on picking up paper
        \item First attempt: incomplete grasp, object slips
        \item Policy automatically re-attempts grasp
        \item Event signal behavior (Fig.~9):
          \begin{itemize}
            \item Remains zero during failed attempt
            \item Rises to 1 only when successful lift occurs
          \end{itemize}
        \item Joint trajectories show multiple grasp attempts
      \end{itemize}
    \end{column}
    \begin{column}{0.45\textwidth}
      \centering
      \includegraphics[width=0.7\textwidth,keepaspectratio]{extracted/figure_9.png}
      \captionof{figure}{Diffusion policy recovering from failed grasp attempts (Fig.~9).}
    \end{column}
  \end{columns}
\end{frame}

%------------------------------------------------
\begin{frame}
  \frametitle{Combined VLA + Diffusion: Task Execution}
  \small
  \begin{columns}
    \begin{column}{0.55\textwidth}
      \begin{itemize}
        \item Full pipeline for pick-and-place (Fig.~10):
          \begin{itemize}
            \item VLA: approach object and plate
            \item Diffusion: grasp and release
            \item Event signal $\sigma$ drives switching
          \end{itemize}
        \item Trajectories:
          \begin{itemize}
            \item Vertical motion and $\sigma$ over time show clear phases
            \item x-y trajectories show distinct VLA vs. diffusion segments
          \end{itemize}
        \item Color coding:
          \begin{itemize}
            \item Purple: VLA-controlled motion
            \item Green: diffusion-controlled motion
          \end{itemize}
      \end{itemize}
    \end{column}
    \begin{column}{0.45\textwidth}
      \centering
      \includegraphics[width=0.7\textwidth,keepaspectratio]{extracted/figure_10.png}
      \captionof{figure}{Combined VLA-diffusion pipeline and success scores (Fig.~10).}
    \end{column}
  \end{columns}
\end{frame}

%------------------------------------------------
\begin{frame}
  \frametitle{Combined VLA + Diffusion: Performance}
  \small
  \begin{itemize}
    \item Success scores for all object-plate combinations:
      \begin{itemize}
        \item VLA + diffusion: $>$ 0.8 for all combinations
        \item VLA-only (trained on pepper and tape only):
          \begin{itemize}
            \item Pepper: 0.43
            \item Tape: 0.20
          \end{itemize}
      \end{itemize}
    \item Hybrid approach:
      \begin{itemize}
        \item More than doubles success compared to VLA-only in many cases
        \item Enables robust manipulation of more challenging objects (e.g., tape, paper)
      \end{itemize}
    \item Demonstrates effective integration of high-level and low-level controllers
  \end{itemize}
\end{frame}

%------------------------------------------------
\begin{frame}
  \frametitle{Conclusion}
  \small
  \begin{itemize}
    \item Introduced a \textbf{hybrid VLA-diffusion framework} for dexterous manipulation
    \item First demonstration (to authors' knowledge) of VLA control on a multi-fingered hand
    \item Key contributions:
      \begin{itemize}
        \item Event-signal-based switching between VLA and diffusion policies
        \item Use of concatenated multi-view images to improve VLA localization
        \item Demonstration of multi-modal grasping and failure recovery via diffusion
      \end{itemize}
    \item Achieved $>$80\% success in language-conditioned pick-and-place
      \begin{itemize}
        \item Significant improvement over VLA-only baseline
      \end{itemize}
  \end{itemize}
\end{frame}

%------------------------------------------------
\begin{frame}
  \frametitle{Limitations and Future Work}
  \small
  \begin{itemize}
    \item \textbf{Limitations:}
      \begin{itemize}
        \item One diffusion policy per object (lookup table selection)
        \item openVLA pre-trained only on 1 DoF pinch grippers
        \item Limited object set and task diversity
      \end{itemize}
    \item \textbf{Future directions:}
      \begin{itemize}
        \item More generalized diffusion policies or use of large-scale datasets
        \item VLA models trained on diverse multi-fingered hands
        \item Incorporation of variable stiffness as a control input
        \item Integration of tactile sensing for improved interaction control
      \end{itemize}
    \item Explore hardware-software co-design to further enhance robustness
  \end{itemize}
\end{frame}

%------------------------------------------------
\begin{frame}
  \centering
  \vfill
  \Large Thank you for your attention!\\[0.5cm]
  \normalsize Questions?
  \vfill
\end{frame}

%------------------------------------------------
\end{document}